\documentclass[11pt]{article}
\usepackage[utf8]{inputenc}
\usepackage[T1]{fontenc}
\usepackage[a4paper,margin=1in]{geometry}
\usepackage{lmodern}
\usepackage{microtype}
\usepackage{amsmath,amssymb,mathtools}
\usepackage{mathrsfs}
\usepackage{tikz-cd}
\usetikzlibrary{arrows}
\usepackage{enumitem}
\usepackage{longtable,booktabs,array}
\usepackage[hidelinks]{hyperref}
\setlength{\parindent}{0pt}
\setlength{\parskip}{0.55em}
\emergencystretch=3em
\hbadness=10000
\hfuzz=20pt
\usepackage[normalem]{ulem}
\newcommand{\K}{K}
\newcommand{\Gr}{\operatorname{Gr}}
\newcommand{\Fil}{\operatorname{Fil}}
\newcommand{\ch}{\operatorname{ch}}
\newcommand{\Todd}{\operatorname{Todd}}
\newcommand{\cl}{\operatorname{cl}}
\newcommand{\Parf}{\operatorname{Parf}}
\newcommand{\Tor}{\operatorname{Tor}}
\newcommand{\Spec}{\operatorname{Spec}}
\newcommand{\RRR}{\mathrm{RRR}}
\providecommand{\codim}{\operatorname{codim}}
\providecommand{\Inv}{\operatorname{Inv}}
\providecommand{\Pic}{\operatorname{Pic}}
\providecommand{\Z}{\mathbb Z}
\providecommand{\Q}{\mathbb Q}
\providecommand{\Ahat}{\widehat A}
\providecommand{\Ahp}{\widehat A^{+}}
\providecommand{\That}{\widehat T}
\providecommand{\rad}{\operatorname{rad}}
\providecommand{\Gl}{\operatorname{Gl}}
\providecommand{\GL}{\operatorname{Gl}}
\providecommand{\ch}{\operatorname{ch}}
\providecommand{\Todd}{\operatorname{Todd}}
\providecommand{\Gr}{\operatorname{Gr}}
\providecommand{\red}{\text{réd}}

% Core macro additions
\providecommand{\Atil}{\widetilde A}
\providecommand{\Ctil}{\widetilde C}
\providecommand{\K}{\mathrm K}
\providecommand{\cC}{\mathcal C}
\providecommand{\cF}{\mathscr F}
\providecommand{\eps}{\varepsilon}
\providecommand{\ellc}{\ell}
\providecommand{\Tor}{\operatorname{Tor}}
\providecommand{\Coker}{\operatorname{Coker}}
\providecommand{\Ker}{\operatorname{Ker}}
\providecommand{\Supp}{\operatorname{Supp}}
\providecommand{\codim}{\operatorname{codim}}
\providecommand{\Gr}{\operatorname{G}}
\providecommand{\rank}{\operatorname{rank}}
\providecommand{\rang}{\operatorname{rang}}
\providecommand{\cO}{\mathcal O}
\providecommand{\sgacircled}[1]{\tikz[baseline=(char.base)]\node[shape=circle,draw,inner sep=.35pt](char){\mbox{\scriptsize #1}};}

\hypersetup{hypertexnames=false}
\newcommand{\sgastandalonefootnotetext}[2]{%
  \begingroup
  \NoHyper
  \footnotetext{#2}%
  \endNoHyper
  \endgroup
}
\providecommand{\R}{\mathbb R}
\providecommand{\C}{\mathbb C}
\providecommand{\N}{\mathbb N}
\providecommand{\Zell}{\mathbb Z_\ell}
\providecommand{\Qell}{\mathbb Q_\ell}
\providecommand{\C}{\mathbb C}
\providecommand{\Ch}{\operatorname{Ch}}
\providecommand{\Sym}{\operatorname{Sym}}
\providecommand{\Td}{\operatorname{Td}}
\begin{document}
\providecommand{\ob}{\operatorname{ob}}
\providecommand{\Ob}{\operatorname{ob}}
\begin{enumerate}[label={[\arabic*]}]
\setcounter{enumi}{13}
\item
% from main.tex
%%%%%%%%%%  scan idx 234  |  folio 221  %%%%%%%%%%
Lang, S., Introduction to differentiable manifolds,
\item \L{}ojasiewicz, S., Ann. della Scuola Norm. Sup. Pisa, serie III, 18, fasc. IV, 1964.
\item Palais, R., Exposés au Séminaire Princeton 1963-64 sur le théorème de l'index.
\item Shih, W., Fiber cobordism and the index of a family of elliptic differential operators,
\item Verdier, J.-L., Les catégories dérivées des catégories abéliennes, à paraître dans North Holl. Pub. Comp.
\end{enumerate}



% Exposé III additions
\renewcommand{\Ob}{\operatorname{Ob}}
\providecommand{\qcoh}{\mathrm{qcoh}}
\providecommand{\coh}{\mathrm{coh}}
\providecommand{\torf}{\mathrm{torf}}
\providecommand{\toramp}{\operatorname{tor.amp}}
\providecommand{\amp}{\operatorname{amp}}
\providecommand{\Id}{\operatorname{Id}}
\providecommand{\pr}{\operatorname{pr}}
\providecommand{\coker}{\operatorname{coker}}
\providecommand{\im}{\operatorname{Im}}
\providecommand{\Spec}{\operatorname{Spec}}



% from main.tex
%%%%%%%%%%  scan idx 235  |  folio 222  %%%%%%%%%%
\section*{\underline{EXPOSÉ III}\\[3pt] \underline{\underline{CONDITIONS DE FINITUDE RELATIVES}}}
\begin{center}
\underline{par L. ILLUSIE}
\end{center}

\section*{1. Pseudo-cohérence relative}

\textbf{Théorème 1.1.} Soit \(n\in\mathbb Z\). Pour chaque morphisme localement de type fini de schémas \(f:X\longrightarrow Y\), il existe une et une seule sous-catégorie pleine \(D(f)_{n\text{-}\coh}\) de \(D(X)\)\footnote{\(D(X)\) = catégorie dérivée de celle des \(\mathcal O_X\)-Modules.} telle que les conditions ci-dessous soient satisfaites :
\begin{enumerate}[label=(\roman*)]
\item Si \(E\in\Ob D(f)_{n\text{-}\coh}\), alors \(i^*(E)\in\Ob D(fi)_{n\text{-}\coh}\) pour tout ouvert \(i:U\to X\).

\item Si \(E\in\ob D(X)\) et s'il existe un recouvrement ouvert \((i_\alpha:U_\alpha\to X)_{\alpha\in A}\) tel que \(i_\alpha^*(E)\in\ob D(fi_\alpha)_{n\text{-}\coh}\) pour tout \(\alpha\in A\), alors \(E\in\ob D(f)_{n\text{-}\coh}\).

\item Si
\[
\begin{tikzcd}
X \arrow[rr,"f"] \arrow[dr,"h"'] & & Y \arrow[dl,"g"]\\
& Z &
\end{tikzcd}
\]
est un triangle commutatif de morphismes localement de type fini, où \(f\) est une immersion fermée et \(g\) un morphisme lisse, alors on a
\[
  E\in\Ob D(h)_{n\text{-}\coh}
  \quad\Longleftrightarrow\quad
  f_*(E)\in D(Y)_{n\text{-}\coh}\qquad\mbox{\footnotemark}\ .
\]
\footnotetext{\(D(Y)_{n\text{-}\coh}\) = sous-catégorie pleine de \(D(Y)\) formée des objets \(n\)-pseudo-cohérents (I~2.15).}
\end{enumerate}

\emph{Démonstration.} L'\emph{unicité} est évidente. Soit en effet \(f:X\to Y\) un morphisme localement de type fini. Choisissons un recouvrement ouvert \((i_\alpha:X_\alpha\to X)_{\alpha\in A}\) tel que, pour chaque \(\alpha\in A\), \(fi_\alpha\) se factorise en
\[
  X_\alpha\xrightarrow{i'_\alpha}Y_\alpha\xrightarrow{f'_\alpha}Y,
\]
où \(i'_\alpha\) est une immersion fermée et \(f'_\alpha\) un morphisme lisse. Il résulte aussitôt de (i), (ii), (iii) que les objets de \(D(f)_{n\text{-}\coh}\) sont les objets \(E\) de \(D(X)\) tels que
\[
  i'_{\alpha *}i_\alpha^*(E)\in \ob D(Y_\alpha)_{n\text{-}\coh};
\]
donc \(D(f)_{n\text{-}\coh}\) est

% from main.tex
%%%%%%%%%%  scan idx 236  |  folio 223  %%%%%%%%%%
bien déterminée (puisque c'est une sous-catégorie pleine de \(D(X)\)). Pour prouver l'existence, nous aurons besoin de quelques lemmes.

\textbf{Lemme 1.1.1.} Soit \(i:X\to Y\) une immersion fermée de schémas telle que \(i_*(\mathcal O_X)\) soit pseudo-cohérent, et soit \(E\in\ob D(X)\). Pour que \(E\) soit \(n\)-pseudo-cohérent, il faut et il suffit que \(i_*(E)\) le soit.

\emph{Preuve.} Le « il faut » résulte du critère général de stabilité de la pseudo-cohérence par application d'un foncteur (I~2.13). Inversement, supposons que \(i_*(E)\) est \(n\)-pseudo-cohérent, et prouvons qu'il en est de même de \(E\). Comme \(i_*(E)\) est à cohomologie localement bornée supérieurement, il en est de même de \(E\), donc, la question étant locale, on peut supposer que \(E\) est \(m\)-pseudo-cohérent pour \(m\) assez grand. Montrons, par récurrence descendante sur \(m\), que \(E\) est \(m\)-pseudo-cohérent pour tout \(m\geq n\). On peut supposer déjà prouvé que \(E\) est \((n+1)\)-pseudo-cohérent. Quitte à localiser, on a donc un triangle distingué
\[
\begin{tikzcd}
& E'' \arrow[dl,"|" marking] & \\
E' \arrow[rr] & & E \arrow[ul]
\end{tikzcd}
\]
où \(E'\) est parfait et \(E''^i=0\) pour \(i>n\). D'après (I~2.6), on est ramené à montrer que \(E''\) est \(n\)-pseudo-cohérent. Or, appliquant \(i_*\) au triangle précédent, on obtient un triangle distingué
\[
\begin{tikzcd}
& i_*(E'') \arrow[dl,"|" marking] & \\
i_*(E') \arrow[rr] & & i_*(E) \arrow[ul]
\end{tikzcd}\ .
\]
D'après la première partie de la démonstration, \(i_*(E')\) est pseudo-cohérent, donc en vertu de (I~2.6), \(i_*(E'')\) est \(n\)-pseudo-cohérent. Comme \(i_*(E'')\) est nul en degré \(>n\), cela signifie, d'après (I~2.10 b)), que \(H^n(i_*(E'')\) est de type fini. Comme \(H^n(i_*(E''))\simeq i_*H^n(E'')\), on est ramené à montrer, d'après (I~2.10 b)), que si \(M\) est un \(\mathcal O_X\)-Module tel que \(i_*(M)\) soit de type fini, alors \(M\) est de type fini, ce qui est immédiat (et indépendant de l'hypothèse sur \(\bar i_*(\mathcal O_X)\)).


% from main.tex
%%%%%%%%%%  scan idx 237  |  folio 224  %%%%%%%%%%
\textbf{Remarque 1.1.2.} Le lemme précédent s'applique notamment au cas où \(i\) est une \emph{immersion régulière} (VII~1.4). En effet, on a localement, grâce au complexe de Kozsul, une résolution gauche finie de \(i_*(\mathcal O_X)\) par des \(\mathcal O_Y\)-Modules libres de type fini, donc \(i_*(\mathcal O_X)\) est parfait (I 4.8 et 4.14) (donc a fortiori pseudo-cohérent).

\textbf{Lemme 1.1.3.} Soit un carré cartésien de schémas
\[
\begin{tikzcd}
X \arrow[d,"i"'] & X' \arrow[l,"f"'] \arrow[d,"i'"]\\
Y & Y'{,} \arrow[l,"g"']
\end{tikzcd}
\]
où \(i\) est une immersion fermée, et soit \(E\) un complexe de \(\mathcal O_X\)-Modules. Alors la flèche canonique de changement de base (EGA~III 1.4.15.4)
\[
  g^*i_*(E)\longrightarrow i'_*f^*(E)
\]
est un isomorphisme.

(ii) Posons \(E'=f^*(E)\). Si \(g\) est plat et \(i_*(E)\) \(n\)-pseudo-cohérent, alors \(i'_*(E')\) est \(n\)-pseudo-cohérent. Inversement, si \(g\) est fidèlement plat et \(i'_*(E')\) \(n\)-pseudo-cohérent, alors \(i_*(E)\) est \(n\)-pseudo-cohérent.

\emph{Preuve.} L'assertion se vérifie trivialement sur les fibres.

\textbf{Lemme 1.1.4.}

% from main.tex
%%%%%%%%%%  scan idx 238  |  folio 225  %%%%%%%%%%
Soit un diagramme commutatif de schémas
\[
\begin{tikzcd}
X' \arrow[dr,"f'"'] & X \arrow[l,hook',"i'"'] \arrow[r,hook,"i''"] & X'' \arrow[dl,"f''"]\\
& Y{,} &
\end{tikzcd}
\]
où \(i'\) (resp. \(i''\)) est une immersion fermée, et \(f'\) (resp. \(f''\)) un morphisme lisse, et soit \(E\in\ob D(X)\). On a alors :
\[
  i'_*(E)\ n\text{-pseudo-cohérent}
  \quad\Longleftrightarrow\quad
  i''_*\ n\text{-pseudo-cohérent}.
\]

\emph{Preuve.} Suivant une astuce due à Lichtenbaum, on envoie \(X\) diagonalement dans \(X'\times_YX''\), et l'on est aussitôt ramené à prouver l'assertion dans le cas où l'un des morphismes \(f'\), \(f''\) est l'identité. Changeant les notations, on a donc un triangle commutatif
\begin{equation}\tag{$*$}
\begin{tikzcd}
& Y \arrow[d,"f"]\\
X \arrow[r,hook,"j"'] \arrow[ur,hook,"i"] & Z{,}
\end{tikzcd}
\end{equation}
où \(i\) et \(j\) sont des immersions fermées et \(f\) un morphisme lisse, et l'on doit montrer que « \(i_*(E)\) \(n\)-pseudo-cohérent » équivaut à « \(j_*(E)\) \(n\)-pseudo-cohérent ». La question étant locale dans \(Y\) et dans \(Z\) au voisinage d'un point de \(X\), on peut supposer que \(X,Y,Z\) sont affines et qu'on a une factorisation
\[
\begin{tikzcd}
Y \arrow[r,hook,"i'"] \arrow[d,"f"'] & Z[T_1,\ldots,T_r]\arrow[dl,"f'"]\\
Z{,}&
\end{tikzcd}
\]
où \(i'\) est une immersion fermée, et la flèche oblique la projection canonique. Alors (VII 1.10) \(i'\) est une immersion régulière, donc (1.1.1) un objet \(F\) de \(D(Y)\) est \(n\)-pseudo-cohérent si et seulement si \(i'_*(F)\) l'est.


\end{document}
